\section{Related Works}
\subsection{Static Neural Rendering}
In recent years, we have witnessed significant progress in the field of novel view synthesis empowered by Neural Radiance Fields. While vanilla NeRF~\cite{NeRF} manages to synthesize photorealistic images for any viewpoint using MLPs, numerous subsequent works focus on acceleration~\cite{Plenoxels, EfficientNeRF, SNeRF, Instant-NGP}, real-time rendering~\cite{MobileNeRF, PlenOctrees}, camera parameter optimization~\cite{NoPe-NeRF, BARF, FF-NeRF}, few-shot learning~\cite{FDNeRFF, FreeNeRF, pixelNeRF}, unbounded scenes~\cite{Mip-NeRF360, Omni-NeRF}, improving visual quality~\cite{Mip-NeRF, Zip-NeRF}, and so on.

More recently, a novel framework 3D Gaussian Splatting~\cite{3D-GS} has received widespread attention for its ability to synthesize high-fidelity images for complex scenes in real-time along with a fast training speed. The key insight lies in that it exploits a point-like representation, referred to as 3D Gaussians. However, these works are mainly restricted to the domain of static scenes.

\subsection{Dynamic Neural Rendering}

To extend neural rendering to dynamic scenes, current efforts primarily focus on deformation-based~\cite{D-NeRF, HyperNeRF, NR-NeRF} and flow-based methods ~\cite{DyNeRF, NSFF, NeRFlow, VideoNeRF}. However, these approaches share similar issues as NeRF, including slow training and rendering speed. To mitigate the efficiency problem, various acceleration techniques like voxel~\cite{TiNeuVox, DeVRF} or hash-encoding representation~\cite{TI-DNeRF}, and spatial decomposition~\cite{KPlanes, Tensor4D, HexPlane, D2NeRF} have emerged. Given the increased complexity brought by dynamic scene modeling, there still exists a gap in rendering quality, training time, and rendering speed compared to static scenes. 

Concurrent with our work, methods like Deformable 3D Gaussians (D-3D-GS)~\cite{Deformable3DGS}, 4D-GS~\cite{RealTime4DGS}, and Dynamic 3D Gaussians~\cite{Dynamic3DGS} leverage 3D-GS as the scene representation, expecting this novel point-like representation can facilitate dynamic scene modeling. D-3D-GS directly integrates a heavy deformation network into 3D-GS, while 4D-GS combines HexPlane~\cite{HexPlane} with 3D-GS to achieve real-time rendering and superior visual quality. Dynamic 3D Gaussians proposes a method that simultaneously addresses the tasks of dynamic scene novel-view synthesis and 6-DOF tracking of all dense scene elements. While our method also takes 3D-GS as the scene representation, unlike any of the aforementioned methods, our main motivation is to aggregate 3D Gaussians with similar deformations into a superpoint to significantly decrease the computational expense required.

\subsection{Superpixel/Superpoint}
There exists a long line of research works on superpixel/superpoint segmentation and we refer readers to the recent paper~\cite{J2023AnES} for a thorough survey. Here we focus on neural network-based methods. 

On one hand, methods including SFCN~\cite{SFCN}, AINet~\cite{AINet}, and LNS-Net~\cite{LNS-Net} adopt a neural network for generating superpixels. SFCN utilizes a fully convolutional network associated with an SLIC loss, while AINet introduces an implantation module and a boundary-perceiving SLIC loss for generating superpixels with more accurate boundaries. LNS-Net proposes an online learning framework, alleviating the demand for large-scale manual labels. On the other hand, existing methods for point cloud over-segmentation can be divided into two categories: optimization-based methods~\cite{Papon2013VoxelCC, Lin2018TowardBB, Guinard2017WEAKLYSS, Landrieu2016CutPF} and deep learning-based methods ~\cite{Landrieu2019PointCO, Hui_2023_ICCV, Hui2021SuperpointNF}.

Our approach can be treated as an over-segmentation of dynamic point clouds, which is an unexplored realm. Existing superpixel/superpoint methods cannot be directly applied to our task since it is challenging to maintain superpoint-segmentation consistency across the temporal domain. Moreover, prevalent methods either employ computationally intensive backbones or do not support parallelization, making the segmentation a heavy module, which will hinder the fast training and rendering speed of our approach.